\documentclass[a4paper,oneside,12pt]{book}

%----------------------------------------------------------------------------------------
%	README!
%   Welcome. It's worth having a read through this file
%   to set up the broad parameters, such as the name of
%   the degree, the school/department, the type of work
%   (dissertation/Final Year Project/report, etc. as well
%   as your own details.
%----------------------------------------------------------------------------------------

%----------------------------------------------------------------------------------------
%	COVER PAGE
%   The cover page is laid out in title/title.tex. You can choose a colour
%   or black and white logo
%----------------------------------------------------------------------------------------

%----------------------------------------------------------------------------------------
%	THESIS INFORMATION
%   Put title, author name, supervisor name, degree, type of work, school, department in here
%   It will be used for the title page and for the embedded PDF information
%----------------------------------------------------------------------------------------

\newcommand{\thesistitle}{``stealthy'' ECH} % Your thesis title, this is used in the title and abstract
\newcommand{\degree}{Msc. Computer Science -- Data Science} % Replace with your degree name, this is used in the title page and abstract
\newcommand{\typeofthesis}{dissertation} % dissertation, Final Year Project, report, etc.
\newcommand{\authorname}{Neimhin Christopher Robinson Gunning} % Your name, this is used in the title page and PDF stuff
%% Do not put your Student ID in the document, as TCD will not publish
%% documents that contain both your name and your Student ID.
\newcommand{\supervisor}{Dr. Stephen Farrell} % replace with the name of your supervisor
%\newcommand{\cosupervisor}{Dr Alex Lee} % replace with the name of your co-supervisor if you have one
\newcommand{\keywords}{censorship, privacy, security, ECH, steganography} % Replace with keywords for your thesis
\newcommand{\school}{\href{https://scss.tcd.ie}{School of Computer Science and Statistics}} % Your school's name and URL, this is used in the title page
%Edited by HS for engineering

%% Comment out the next line if you don't want a department to appear
%\newcommand{\department}{\href{http://researchgroup.university.com}{Department Name}} % Your research group's name and URL, this is used in the title page


%% Language and font encodings
\usepackage[T1]{fontenc} 
\usepackage[utf8]{inputenc}
\usepackage[english]{babel}

%% Bibliographical stuff
\usepackage[]{cite}
%% Document size
% include showframe as an option if you want to see the boxes
\usepackage[a4paper,top=2.54cm,bottom=2.54cm,left=2.54cm,right=2.54cm,headheight=16pt]{geometry}
\setlength{\marginparwidth}{2cm}
%% Useful packages
\usepackage{amsmath}
\usepackage[autostyle=true]{csquotes} % Required to generate language-dependent quotes in the bibliography
\usepackage[pdftex]{graphicx}
\usepackage[colorinlistoftodos]{todonotes}
\usepackage[colorlinks=true, allcolors=black]{hyperref}
\usepackage{hyperxmp}
\usepackage{caption} % if no caption, no colon
\usepackage{sfmath} %use sans-serif in the maths sections too
\usepackage[parfill]{parskip}    % Begin paragraphs with an empty line rather than an indent
\usepackage{setspace} % to permit one-and-a-half or double spacing
\usepackage{enumerate} % fancy enumerations like (i) (ii) or (a) (b) and suchlike
\usepackage{booktabs} % To thicken table lines
\usepackage{fancyhdr}
\usepackage{xcolor} % to get TCD colour on headings
\usepackage{todonotes}
\numberwithin{equation}{chapter} %HS edit for (chapter.equation)
\pagestyle{plain} % Embrace simplicity!

\definecolor{tcd_blue}{RGB}{5, 105, 185}

%% It's personal taste but...
%% Uncomment the following block if you want your name and ID at the top of
%% (almost) every page.

%\pagestyle{fancy}
%\fancyhf{} % sets both header and footer to nothing
%\renewcommand{\headrulewidth}{0pt}
%\cfoot{\thepage}
%\ifdefined\authorid
%\chead{\it \authorname\ (\authorid)}
%\else
%\chead{\it \authorname}
%\fi
%% End of block

%% It is good practise to make your font sans-serif to improve the accessibility of your document.  Comment out the following line to disable it (but you really should not)
\renewcommand{\familydefault}{\sfdefault} %use the sans-serif font as default

%% If you insist on not using sans-serif (please don't), consider using Palatino instead of the LaTeX standard
%\usepackage{mathpazo} % Use the Palatino font by default if you prefer it to Computer Modern


%% Format Chapter headings appropriately
\usepackage{titlesec}
\titleformat{\chapter}[hang]{\normalfont\huge\bfseries\color{tcd_blue}}{\thechapter}{1cm}{}{}

\title{\thesistitle}
\author{\authorname}


\hypersetup{
   pdftitle=\thesistitle, % Set the PDF's title to your title
   pdfauthor=\authorname, % Set the PDF's author to your name
   pdfkeywords=\keywords, % Set the PDF's keywords to your keywords
   pdfsubject=\degree, % Set the PDF's keywords to your keywords
   pdfinfo={
     pdfsupervisor=\supervisor, % Set the PDF's supervisor to your supervisor
     %pdfcosupervisor=\cosupervisor, % Set the PDF's cosupervisor to your cosupervisor if using
   }
}


\frontmatter
\begin{document}
\begin{titlepage}

\center % Center everything on the page

%% All the text parameters should be taken from the start of the main.tex file.
%% You should only alter stuff here if you want to change the layout

%----------------------------------------------------------------------------------------
%	LOGO SECTION
%----------------------------------------------------------------------------------------
%% Choose one of the following -- a colour or black-and-white logo

\includegraphics{title/Trinity_RGB_transparent_main.png}\\[1cm] 
%\includegraphics[width=12cm]{title/black-stacked-trinity.jpg}\\[1cm] 

\Large \school\\[1.5cm] % Minor heading such as course title
\ifdefined\department
\large \department\\[1.5cm] % Minor heading such as course title
\fi

%----------------------------------------------------------------------------------------
%	TITLE SECTION
%----------------------------------------------------------------------------------------
\makeatletter
{ \huge \bfseries \thesistitle}\\[1.5cm] % Title of your document
 

%----------------------------------------------------------------------------------------
%	AUTHOR SECTION
%----------------------------------------------------------------------------------------

\ifdefined\authorid
\authorname\\ % Your name
\authorid\\[2cm] % Your Student ID
\else
\authorname\\[2cm] % Your name
\fi

%----------------------------------------------------------------------------------------
%	Supervisor SECTION
%----------------------------------------------------------------------------------------

Supervisor: \supervisor\\[2cm] % Their name
\ifdefined\cosupervisor
Cosupervisor: \cosupervisor\\[2cm] % Their name
\fi


%----------------------------------------------------------------------------------------
%	DATE SECTION
%----------------------------------------------------------------------------------------

{\large \today}\\[2cm] % Date, change the \today to a set date if you want to be precise

 
%----------------------------------------------------------------------------------------
%	TYPE OF THESIS SECTION
%----------------------------------------------------------------------------------------
\vfill
 A \typeofthesis\ submitted in partial fulfilment\\of the requirements for the degree of\\
\degree

\vfill % Fill the rest of the page with whitespace

\end{titlepage}

\pagenumbering{roman}
\section*{\Huge\textcolor{tcd_blue}{Declaration}}
\vspace{1cm}
I hereby declare that this \typeofthesis\ is entirely my own work and that it has not been submitted as an exercise for a degree at this or any other university.

\vspace{1cm}
I have read and I understand the plagiarism provisions in the General Regulations of the University Calendar for the current year, found at \url{http://www.tcd.ie/calendar}.
\vspace{1cm}

I have completed the Online Tutorial on avoiding plagiarism `Ready Steady Write', located at \url{http://tcd-ie.libguides.com/plagiarism/ready-steady-write}.
\vspace{1cm}

I consent to the examiner retaining a copy of the thesis beyond the examining period, should they so wish (EU GDPR May 2018).
\vspace{1cm}

I agree that this thesis will not be publicly available, but will be available to TCD staff and students in the University’s open access institutional repository on the Trinity domain only, subject to Irish Copyright Legislation and Trinity College Library conditions of use and acknowledgement.
\todo{
  \textbf{Please consult with your supervisor on this last item before agreeing, and delete if you do not consent}
}
\vspace{3cm}

Signed:~\rule{5cm}{0.3pt}\hfill Date:~\rule{5cm}{0.3pt}

\chapter*{Abstract}
The TLS protocol does not encrypt the Server Name Indication (SNI) nor the Application Layer Protocol Negotiation (ALPN) lists from IP nodes facilitating a TLS connection.
These data are sensitive and it is desirable in many cases to keep them private. Encrypted Client Hello (ECH) is a proposed solution which encrpyts the entire ClientHello message in which
the SNI and ALPN are contained, however ECH can only be carried out with the agreement of all facilitating nodes. This work proposes TODO how many? a suite of possible designs for a more ``stealthy'' version of
strategy for hiding the sensitive parts of the CLientHello. These strategies should not alert intermediate nodes that the SNI and ALPN list is being kept secret.
Evalutions of the implementation of a subset of these designs are presented.

\chapter*{Lay Abstract}
Similar to the actual abstract in terms of the information, but written for a non-specialist. So no jargon, no acronyms. Explain to a member of the general public what this project entailed. Should be no longer than the actual abstract.


\todo{
  This must be on a separate page.
}

\newpage
\onehalfspacing\raggedright %\raggedright turns off justification and hyphenation

\section*{\Huge\textcolor{tcd_blue}{Acknowledgements}}
\todo {
  Thanks Everyone!

}
  You should acknowledge any help that you have received (for example from technical staff), or input provided by, for example, a company.
\tableofcontents
\listoffigures
\listoftables
\newpage
\section*{\Huge\textcolor{tcd_blue}{Nomenclature}}
\begin{tabular}{lp{9cm}l}
A&Area of the wing&$m^{2}$\\
B\\
C& Roman letters first, with capitals\ldots\\
a&then lower case.\\
b\\
c\\
$\Gamma$&Followed by Greek capitals\ldots\\
$\alpha$&then lower case greek symbols.\\
$\beta$\\
$\epsilon$\\
ECH&Encrypted Client Hello\\
ALPN&Application Layer Protocol Negotiation\\
SNI&Server Name Indication\\
\todo{
  TLA&Finally, three letter acronyms and other abbreviations arranged alphabetically\\
}
\end{tabular}
\vspace{2cm}

If a parameter has a typical unit that is used throughout your report, then it should be included here on the right hand side.

If you have a very mathematical report, then you may wish to divide the nomenclature list into functions and variables, and then sub- and super-scripts.

If you have a large number of acronyms, check out \href{https://www.overleaf.com/learn/latex/Glossaries} to make that more robust.

Note that Roman mathematical symbols are typically in a serif font in italics.

\mainmatter
% maintaining separate .tex files for each chapter is good practice
% \input{introduction/introduction.tex}
% \input{background/background.tex}
% \input{project/project.tex}
% \input{evaluation/evaluation.tex}
% \input{conclusion/conclusion.tex}

% note that your supervisor may have a strong opinion on the style of referencing you use. Some background is available at https://www.overleaf.com/learn/latex/Bibtex_bibliography_styles
\bibliographystyle{IEEEtran} %Changed to IEEETran by HS
%\bibliographystyle{unsrt}

\bibliography{bibs/sample}
\appendix
\renewcommand{\thechapter}{A\arabic{chapter}}
\chapter{Appendix}
You may use appendices to include relevant background information, such as calibration certificates, derivations of key equations or presentation of a particular data reduction method. You should not use the appendices to dump large amounts of additional results or data which are not properly discussed. If these results are really relevant, then they should appear in the main body of the report.

\section{Appendix numbering}
Appendices are numbered sequentially, A1, A2, A3\ldots The sections, figures and tables within appendices are numbered in the same way as in the main text. For example, the first figure in Appendix A1 would be Figure A1.1. Equations continue the numbering from the main text.




\end{document}
